\section{Fr\"olicher, Bott--Chern, Aeppli, and Hermitian metrics}
\label{sec:partC}

\subsection{C.1. The differentials and their dependence on the Hodge number}

In this section $X$ is the standard marked $(3,4)$ threefold and
$h=h^{1,2}(X)=1$ by Part A.  We retain $h$ in the formal calculations
to specify exactly which conclusions depend on its evaluated value.
Every assertion involving $h$ is valid for every value
consistent with the Hodge table
\begin{equation}\label{eq:C-hodge-input}
\begin{array}{c|rrrr}
q\backslash p&0&1&2&3\\\hline
3&0&0&0&1\\
2&0&h&h+1&1\\
1&1&h+1&h&0\\
0&1&0&0&0.
\end{array}
\end{equation}
Thus no conclusion below silently uses $h=1$.  The geometric inputs for
this part are the structure-sheaf calculation, the vanishing of
holomorphic forms, and the sphere topology, which give
\eqref{eq:C-hodge-input}; Statements A1.3 and A3.1 determine $h$.
We use $d=\partial+\bar\partial$ on complex-valued
smooth forms and the complex orientation for integration.

\begin{claim}{C1.1}{Proved}{The complete Fr\"olicher differential in terms of $h$}
For a compact complex threefold with \eqref{eq:C-hodge-input} and the
complex de Rham cohomology of $S^6$, the Fr\"olicher spectral sequence
satisfies $E_1\ne E_2=E_\infty$.  Its only nonzero $d_1$ ranks are
\[
\rank d_1^{0,1}=1,\quad \rank d_1^{1,1}=h,\quad
\rank d_1^{1,2}=h,\quad \rank d_1^{2,2}=1.
\]
In particular its two middle rows are exact.
\end{claim}
\begin{proof}
Here $E_1^{p,q}=H^q(X,\Omega_X^p)$ and
$d_r:E_r^{p,q}\longrightarrow E_r^{p+r,q-r+1}$ are complex-linear.
For $r=2$, any source in the two middle rows has target in one of
$(2,0),(3,0),(3,1)$ or outside the table, except the source $(0,2)$,
which is zero.  Sources in row $q=3$ outside $(3,3)$ are zero, and
$(3,3)$ has no target.  For $r=3$ the only possible targets inside the
table have sources $(0,2)$ or $(0,3)$, both zero.  For $r>3$ there is
no arrow inside the table.  A zero term remains zero on subsequent
pages.  Hence $E_2=E_\infty$.

The filtration on $H^k_{\mathrm{dR}}(X;\C)$ has no nonzero associated
graded term for $1\le k\le5$.  Every middle term of $E_2$ therefore
vanishes.  Consequently
\[
0\longrightarrow\C\longrightarrow\C^{h+1}
\longrightarrow\C^h\longrightarrow0,
\qquad
0\longrightarrow\C^h\longrightarrow\C^{h+1}
\longrightarrow\C\longrightarrow0
\]
are exact, giving the stated ranks.  The first rank is one, so $E_1$
cannot equal $E_2$.
\end{proof}

\begin{claim}{C1.2}{Proved}{A basis on which all first differentials are specified}
Choose a nonzero $g\in H^{0,1}_{\bar\partial}(X)$, a basis
$q_1,\ldots,q_h$ of a complement to $\C d_1g$ in
$H^{1,1}_{\bar\partial}(X)$, a basis $v_1,\ldots,v_h$ of
$H^{1,2}_{\bar\partial}(X)$, and a class $e$ complementary to the
image of $d_1:H^{1,2}\to H^{2,2}$.  With
$a=d_1g$, $c_j=d_1q_j$, $d_j=d_1v_j$, and $t=d_1e$, the first
differential is
\[
g\mapsto a\mapsto0,\quad q_j\mapsto c_j\mapsto0,
\qquad v_j\mapsto d_j\mapsto0,\quad e\mapsto t\mapsto0.
\]
These are bases in the corresponding bidegrees.  If $h=1$ and a
class $q$ restricts to the invariant fibre class $\eta$ described below,
we may take $q_1=q$, $v_1=gq$, and $e=q^2$, after scaling $g$ if
necessary.  Statements A1.3 and A3.2 prove these hypotheses for $X$,
so this geometric choice applies to the threefold in the question.
\end{claim}
\begin{proof}
The first assertions follow from the injectivity and surjectivity
proved in C1.1.  For the last assertion write
$\eta=u\wedge w+6\gamma\wedge\delta$ in the marked integral
cohomology of a regular fibre.  Its complex-oriented square has
integral $-12$.  Its product with the $(0,1)$ image of $\gamma$ is
nonzero: the explicit period calculation in C3.7 gives its nonzero
pairing with $d\zeta_2$.  Thus $gq$ and $q^2$ have nonzero fibre
restrictions.  The images under $d_1$ restrict to zero on the fibre,
because the fibre is a complex torus and its invariant representatives
are $d$-closed.  Therefore $q$ complements $d_1g$, $gq$ spans the
one-dimensional $H^{1,2}$, and $q^2$ complements $d_1H^{1,2}$.
The extension used in this last choice is the actual Leray lift in
Statement A3.2.  This argument uses that extension; it does not infer
it from a calculation on the smooth fibre alone.
\end{proof}

\subsection{C.2. A computation from the two definitions}

\begin{claim}{C2.1}{Proved}{Bott--Chern and Aeppli tables as functions of $h$}
With
\[
H^{p,q}_{BC}=
\frac{\ker\partial\cap\ker\bar\partial\subset A^{p,q}}
{\im\partial\bar\partial},\qquad
H^{p,q}_{A}=
\frac{\ker\partial\bar\partial\subset A^{p,q}}
{\im\partial+\im\bar\partial},
\]
the complete tables are
\begin{equation}\label{eq:C-BC-A}
\begin{array}{c|rrrr}
q\backslash p&0&1&2&3\\\hline
3&0&0&1&1\\
2&0&h&2h&1\\
1&0&2&h&0\\
0&1&0&0&0
\end{array}_{BC}
\qquad
\begin{array}{c|rrrr}
q\backslash p&0&1&2&3\\\hline
3&0&0&0&1\\
2&0&h&2&0\\
1&1&2h&h&0\\
0&1&1&0&0.
\end{array}_{A}
\end{equation}
In particular, substitution of the value $h=1$ proved in A3.1 gives
the workbench's numerical tables.  The totals are
$(1,0,2,2,2,2,1)$ and $(1,2,2,2,2,0,1)$, respectively.
\end{claim}
\begin{proof}
We apply the structure theorem for bounded double complexes over a
field, \citet[Theorem 3]{Stelzig2021}, to the actual algebraic double
complex $(A^{\bullet,\bullet}(X),\partial,\bar\partial)$.
The word bounded refers to the bidegrees, not to the dimensions of the
spaces of smooth forms, so that theorem applies here.  It gives a
direct-sum decomposition by isomorphisms of double complexes into
squares and zigzags.  An odd zigzag has one-dimensional total
cohomology.  All odd zigzags in intermediate total degree are therefore
excluded by $b_1=\cdots=b_5=0$.  An odd zigzag of length greater than
one would have at least two vertices of its cohomological total degree.
The extreme diagonals $p+q=0$ and $p+q=6$ contain just one bidegree
each, so the remaining odd zigzags are exactly one dot at $(0,0)$ and
one at $(3,3)$.

An even zigzag of length $2r$ contributes a nonzero $d_r$ to one of
the two Fr\"olicher spectral sequences.  C1.1 and complex conjugation
exclude $r\ge2$ in both sequences.  Thus the remaining summands are
arrows.  The horizontal arrows, with their multiplicities, are
\begin{equation}\label{eq:C-arrows}
\begin{array}{c|c|r}
\text{source}&\text{target}&\text{multiplicity}\\\hline
(0,1)&(1,1)&1\\
(1,1)&(2,1)&h\\
(1,2)&(2,2)&h\\
(2,2)&(3,2)&1.
\end{array}
\end{equation}
The vertical arrows are their reflections in $p=q$, with the same
multiplicities.  The correspondence with spectral-sequence ranks is
\citet[Proposition 6(1)]{Stelzig2021}; alternatively one computes
column cohomology on each horizontal arrow directly.

We now evaluate the displayed definitions, rather than use a formula
for the dimensions.  On an arrow, the source is killed by neither
the outgoing differential nor its total differential, whereas its
target is closed under both.  There is no composite
$\partial\bar\partial$.  Hence the target contributes one dimension
to Bott--Chern cohomology, and the source contributes one dimension
to Aeppli cohomology.  On a square, only the top right corner is closed
under both operators, and it is the image of
$\partial\bar\partial$ from the bottom left corner.  The other
three corners of the square are images of $\partial$ or
$\bar\partial$, while the bottom left corner is not in
$\ker\partial\bar\partial$.  Thus a square contributes to neither
cohomology.  A dot contributes once to each.  These computations
commute with direct sums, including the possibly infinite sum of
squares.  Counting the sources and targets in \eqref{eq:C-arrows}
and its reflection gives \eqref{eq:C-BC-A}.

The program \texttt{code/c\_double\_complex.py} performs a separate
kernel-and-image computation over $\Q$ and a source-and-target
count.  It includes an acyclic square with its signs, checks
$\partial\bar\partial+\bar\partial\partial=0$, and verifies the
two answers for $h=0,1,2,3,4,5$.  Its output is
\texttt{outputs/c\_double\_complex.txt}.  This is exact linear
algebra, not a floating-point check and not a determination of $h$.
\end{proof}

\begin{claim}{C2.2}{Proved}{Comparison with McHugh and Angella}
The specialization of McHugh's Theorem 6.3, Proposition 6.4, equation
(7), and Section 7.1 to \eqref{eq:C-hodge-input} agrees with C2.1.
The corresponding relations and Figure 5 in Angella's Section 4 also
agree.  No correction to these specific formulas is required.
\end{claim}
\begin{proof}
McHugh's quotient $K^{p,q}$ is zero in intermediate total degrees
on every arrow and square just used: a closed corner of an arrow is
the image of the incoming differential and a closed corner of a square
is in both images.  Thus $k^{p,q}=0$ for $0<p+q<6$.
Also $E_2^{3,1}=0$ by C1.1 and $H^{1,1}(X;\R)=0$ because
$H^2(X;\R)=0$.  Reading the formulas in
\citet[Theorem 6.3, Proposition 6.4, (7), pp. 20--25]{McHugh2018}
therefore gives
\[
h_{BC}^{1,1}=2h^{0,1}=2,\qquad
h_{BC}^{1,2}=h^{1,2}=h,\qquad
h_{BC}^{2,2}=-1+(h+1)+h=2h.
\]
The remaining entries in his table are precisely the boundary entries
of \eqref{eq:C-BC-A}.  In
\citet[Section 4, pp. 13--16]{Angella2018}, the relation
\[
h_{BC}^{2,2}=2h_{BC}^{2,1}-2h^{0,1}+2
\]
gives $2h$, and the boundary values reduce to those above because
$h^{2,0}=h^{0,2}=0$.  This comparison concerns exactly the formulas
used here; it does not certify every assertion in either article.
\end{proof}

\subsection{C.3. The algebra, line bundles, and metric consequences}

\begin{claim}{C3.1}{Proved}{The real Bott--Chern Chern-class map and a global $G$}
Let $G$ be a $\bar\partial$-closed smooth $(0,1)$ representative of
a nonzero $g\in H^1(X,\OO_X)$.  Set
\[
\mathsf A=[\partial G]_{BC},\qquad
\mathsf B=[\bar\partial\bar G]_{BC}.
\]
These classes form a complex basis of $H^{1,1}_{BC}(X)$.
With the exponential $z\mapsto e^{2\pi iz}$ and the Chern convention
$c_{1,BC}=[iF/(2\pi)]$, the map of real Lie groups is
\begin{equation}\label{eq:C-Pic-map}
\Pic(X)=\{\exp(zg):z\in\C\}\longrightarrow H^{1,1}_{BC}(X),
\qquad \exp(zg)\longmapsto z\mathsf A+\bar z\mathsf B.
\end{equation}
It is injective and its image is the real subspace fixed by conjugation.

There is an operator-explicit global representative.  Normalize the
Fubini--Study form by $\int_{\Pone}\omega_{FS}=1$, put
$\theta=f^*\omega_{FS}$, choose any smooth Riemannian metric on $X$,
and let $\mathcal G_d$ be its Green operator for the real de Rham
Laplacian.  Then
\begin{equation}\label{eq:C-G-Green}
\alpha=d^*\mathcal G_d\theta,\qquad G_0=\alpha^{0,1}
\end{equation}
satisfy $d\alpha=\theta$, $\bar\partial G_0=0$,
$[G_0]\ne0$, and $\exp([G_0])=f^*\OO_{\Pone}(1)$.
\end{claim}
\begin{proof}
C1.1 shows that $[\partial G]_{\bar\partial}=d_1g\ne0$.
The conjugate class $\mathsf B$ is nonzero, and its Dolbeault image
is zero.  Thus $\mathsf A,\mathsf B$ are independent.  C2.1 gives
their completeness as a basis.

The exponential sequence and $H^1(X;\Z)=H^2(X;\Z)=0$ identify
$\Pic(X)$ with $H^1(X,\OO_X)=\C g$, with no period lattice.
On the smooth trivial line with unit frame $s$, define
$\bar\partial_zs=-2\pi izG\otimes s$.  This is integrable.
If $G=\bar\partial\phi_i$ locally, the holomorphic frames
$e_i=e^{2\pi iz\phi_i}s$ have transition functions
$e_j=e^{2\pi iz(\phi_j-\phi_i)}e_i$, which represent
$\exp(zg)$.  The metric-compatible Chern connection is
\[
\nabla_z=d-2\pi i(zG+\bar z\bar G),\qquad
F_z=-2\pi i(z\partial G+\bar z\bar\partial\bar G).
\]
This proves \eqref{eq:C-Pic-map}, including its sign.  The basis
proves injectivity and the assertion about its real image.

The form $\theta$ is closed and exact because $H^2(X;\R)=0$.
Hodge decomposition, together with $d\theta=0$ and the commutation
of $d$ with the Green operator, gives $dd^*\mathcal G_d\theta=
\theta$.  The operator preserves real forms.  Since $d\alpha$ has
type $(1,1)$, its $(0,2)$ component says
$\bar\partial G_0=0$.  If $G_0=\bar\partial u$, then
$\theta=\partial\bar\partial(u-\bar u)$.  A real multiple of
$\operatorname{Im}u$ would be a global plurisubharmonic function,
because $\theta$ is semipositive.  The maximum principle on compact
$X$ makes this function constant.  This contradicts the fact that
$\theta$ is nonzero over the regular locus of $f$.

Finally \eqref{eq:C-Pic-map} gives
$c_{1,BC}(\exp([G_0]))=[\theta]=c_{1,BC}(f^*\OO(1))$.
Injectivity gives the asserted line-bundle equality.  Formula
\eqref{eq:C-G-Green} is a global formula involving a specified
elliptic inverse; it is not a finite coordinate formula on the toric
and cyclic charts, and it depends on the auxiliary metric.
\end{proof}

\begin{claim}{C3.2}{Proved}{The factor $12$ with a fixed collar orientation}
Use the marked cohomology class $\gamma$ dual to
$\hat\gamma$, the clockwise meridians and the lift relation of
[AC26, Section 7.5], and put
$p=12l_0-4l_1-3l_2=\pm1$.  Let $g_\gamma$ be the unique coherent
Leray class whose fibre image is $\gamma^{0,1}$.  Then
\begin{equation}\label{eq:C-factor12}
\exp(12g_\gamma)=f^*\OO_{\Pone}(-p).
\end{equation}
For the standard data $(l_0,l_1,l_2)=(0,1,-1)$, this is
$\exp(12g_\gamma)=f^*\OO(1)$.  For the opposite value $p=1$,
the exponent representing $f^*\OO(1)$ is $-12g_\gamma$.
\end{claim}
\begin{proof}
We include the comparison that fixes the sign.  A degree-one integral
class $k\gamma$ on a regular fibre extends over the cyclic filling of
order $m_j$ precisely when the affine generator $\hat g_j$ can be
assigned an integer value $n_j$ satisfying
$m_j n_j=k\gamma(v_j)=kl_j$.  Admissibility gives
$\gcd(l_j,m_j)=1$, so this is precisely $m_j\mid k$.
At the cusp $\gamma$ extends, as follows from its actual specialization
to $H^1(W;\Z)$ in [AC26, Theorem B.1 and Proposition 7.12].
The global monodromy invariants have rank one generated by $\gamma$.
It follows, by intersection of these three local images in the marked
local system, that
$H^0(B,R^1f_*\Z)=12\Z\gamma$.  This recovers the degree-one
part of [AC26, Proposition 7.26], without identifying the reduced and
non-reduced fibre: this is topological cohomology, unchanged by
nilpotents, whereas the coherent comparison uses the full fibre.

Let $U=\bigsqcup_iD_i$ be small discs around the special points and
let $V$ be their pair-of-pants complement with collars.  Lift
$12\gamma$ to integral classes $\Theta_i$ on $f^{-1}(D_i)$ and
$\Theta_V$ on $f^{-1}(V)$.  Existence over $V$ follows from
$H^2(V;\Z)=0$ and the low-degree Leray sequence; over each disc it
follows from the local images just computed.  Write
$a_j=\Theta_V(\hat\rho_j)$.  On the two finite clockwise meridians
the differences $\Theta_i-\Theta_V$ have periods
\[
n_1=4l_1-a_1,\qquad n_2=3l_2-a_2.
\]
At the cusp, the specified lift relation
$\widetilde a_0=t_\mu\hat\rho_2^{-1}\hat\rho_1^{-1}$,
with $\gamma(\mu)=l_0$, and the null-homotopy of
$\widetilde a_0$ in the cusp filling give
$n_0=-12l_0+a_1+a_2$.  Hence
\begin{equation}\label{eq:C-collar-sum}
n_0+n_1+n_2=-p.
\end{equation}
These are the explicitly oriented boundary evaluations in [AC26,
Proposition 7.27, pp. 62--63], not an unspecified sign of a spectral
sequence differential.

Represent the lifts by real closed one-forms.  On a collar their
difference is the pullback of a real closed base form $b_i$ with
clockwise period $n_i$, modulo an exact form.  This follows from the
degree-one Leray sequence over the annulus, since the difference has
zero fibre restriction and the annulus has no second cohomology.
Subtract exact terms on smaller collars, extending their primitives
with cutoffs, so the equality of forms holds wherever a partition of
unity changes.  For a real base cutoff $\rho_i$ equal to one toward
$D_i$ and zero toward $V$, the resulting glued real one-form $\alpha$
on $X$ has
\[
d\alpha=f^*\vartheta,\qquad
\vartheta=\sum_i d\rho_i\wedge b_i.
\]
The form $\vartheta$ is a smooth real base two-form; thus
$\bar\partial\alpha^{0,1}=0$.  Its fibre restriction is the
$(0,1)$ image of $12\gamma$, so the coherent Leray edge isomorphism
$H^1(X,\OO_X)\simeq H^0(B,R^1f_*\OO_X)$ identifies
$[\alpha^{0,1}]=12g_\gamma$.

In a positively oriented polar coordinate $(r,\varphi)$ on $D_i$,
$b_i=-n_i\,d\varphi/(2\pi)$ modulo an exact form.
Because $\rho_i$ decreases from one to zero as $r$ increases,
\[
\int_{D_i\cap V}d\rho_i\wedge b_i
=-n_i\int_0^1d\rho_i=n_i.
\]
The exact part does not change this integral, since $d\rho_i$ is
supported inside the collar.  Therefore
$\int_B\vartheta=-p$ by \eqref{eq:C-collar-sum}.  C3.1 identifies
the Bott--Chern class of $\exp(12g_\gamma)$ with $[f^*\vartheta]$.
On $\Pone$ two real $(1,1)$ forms of the same integral differ by
$i\partial\bar\partial$ of a real function (solve the ordinary
scalar Laplace equation).  Hence this class is the Chern class of
$f^*\OO(-p)$, and injectivity in C3.1 proves
\eqref{eq:C-factor12}.  The program
\texttt{code/c\_sign\_and\_fibre.py} verifies the integer
intersection, \eqref{eq:C-collar-sum}, and the collar integral exactly.
\end{proof}

\begin{claim}{C3.3}{Proved}{The complete Bott--Chern algebra, for every $h$}
The positive-total-degree ideal $N\subset H^{\bullet,\bullet}_{BC}(X)$
satisfies
\[
H^{\bullet,\bullet}_{BC}(X)=\C\cdot1\oplus N,\qquad N^2=0.
\]
Its nonzero bidegrees and dimensions are those of C2.1.  Thus the
square-zero conclusion is independent of the disputed value of $h$.
\end{claim}
\begin{proof}
Take global $\bar\partial$-closed representatives $G,Q_j,V_j,E$
of the classes $g,q_j,v_j,e$ of C1.2.  By the arrow computation in
C2.1, bases for the positive Bott--Chern groups consist of
\[
\partial G,\ \bar\partial\bar G;\qquad
\partial Q_j,\ \bar\partial\bar Q_j;\qquad
\partial V_j,\ \bar\partial\bar V_j;\qquad
\partial E,\ \bar\partial\bar E;\qquad \nu,
\]
where the pairs have bidegrees $(1,1)$, $(2,1),(1,2)$,
$(2,2)$, $(3,2),(2,3)$, respectively, and
$\int_X\nu=1$.  In particular the natural map
$H^{3,2}_{BC}\to H^{3,2}_{\bar\partial}$ is an isomorphism:
on its unique horizontal arrow the closed target maps nontrivially
to column cohomology.

Write $a=\partial G$, $b=\bar\partial\bar G$ and
$c_j=\partial Q_j$.  The graded Leibniz rule gives
\[
\partial\bar\partial(G\wedge\bar G)=-a\wedge b,
\quad
\partial\bar\partial(G\wedge\bar Q_j)
=-a\wedge\bar c_j,
\quad
\partial\bar\partial(\bar G\wedge Q_j)=b\wedge c_j.
\]
Consequently $\mathsf A\mathsf B=0$,
$\mathsf A\overline{\mathsf C_j}=0$, and
$\mathsf B\mathsf C_j=0$.

Let $f^*\OO(1)=\exp(\lambda g)$, where $\lambda\ne0$.
Its curvature may be chosen to be the pullback of a two-form on $B$,
whose square vanishes pointwise.  By C3.1,
$\lambda^2\mathsf A^2+\bar\lambda^2\mathsf B^2=0$.
The Dolbeault image of $\mathsf B^2$ is zero.  Hence
$[a^2]_{\bar\partial}=0$.  But
$a^2=\partial(G\wedge a)$, and C1.1 says that
$d_1:H^{1,2}\to H^{2,2}$ is injective.  Thus
$[G\wedge a]_{\bar\partial}=0$.  Choose a smooth $(1,1)$ form
$\Phi$ satisfying $G\wedge a=\bar\partial\Phi$.  Applying
$\partial$ proves $a^2=\partial\bar\partial\Phi$.
Conjugation gives $\mathsf B^2=0$ as well.

To treat $\mathsf A\mathsf C_j$, pair its Dolbeault image with the
generator $g\in H^{0,1}_{\bar\partial}$.  Since
$\bar\partial c_j=0$ and $G\wedge a=\bar\partial\Phi$,
\[
\int_X a\wedge c_j\wedge G
=-\int_X\bar\partial(\Phi\wedge c_j)=0.
\]
Dolbeault--Serre duality and the isomorphism
$H^{3,2}_{BC}\to H^{3,2}_{\bar\partial}$ give
$\mathsf A\mathsf C_j=0$.  Complex conjugation gives
$\mathsf B\overline{\mathsf C_j}=0$.

Every remaining product whose bidegree lies in $[0,3]^2$ lands in
$(3,3)$: it is a product of a $(1,1)$ class with a $(2,2)$ class,
or of a $(2,1)$ class with a $(1,2)$ class.  Each non-top generator
displayed above is $d$-exact, since its chosen primitive is
$\bar\partial$-closed, or is the conjugate of such a form.
The products therefore have integral zero.  Integration is an
isomorphism $H^{3,3}_{BC}\to\C$ by C2.1 and the class of $\nu$,
so all these products vanish.  All other products exceed the complex
dimension in at least one bidegree.  This exhausts the possibilities.
\end{proof}

\begin{claim}{C3.4}{Proved}{Non-K\"ahler degrees and the precise degreewise statement}
For $k=0,\ldots,6$ the non-K\"ahler degrees are
\begin{equation}\label{eq:C-delta}
(\Delta^k)=(0,2,2+2h,4h,2+2h,2,0).
\end{equation}
If $h\ge1$, the two maps
$\bigoplus_{p+q=k}H^{p,q}_{BC}\to H^k_{\mathrm{dR}}(X;\C)
\to\bigoplus_{p+q=k}H^{p,q}_A$
are both isomorphisms exactly for $k=0,6$.
For the weaker assertion that a $d$-exact pure $d$-closed form of
degree $k$ is $\partial\bar\partial$-exact, the valid degrees are
$k=0,1,6$, and also $k=3,4$ if $h=0$.
\end{claim}
\begin{proof}
The totals from C2.1 are
\[
(h^k_{BC})=(1,0,2,2h,2h,2,1),\qquad
(h^k_A)=(1,2,2h,2h,2,0,1).
\]
Subtracting twice the Betti numbers gives \eqref{eq:C-delta}.
On the two endpoint dots both comparison maps are isomorphisms;
all intermediate de Rham groups vanish.  The first assertion follows
by inspecting when both outside groups vanish.  The weaker assertion
is exactly injectivity of the first comparison map.  In intermediate
degree this is equivalent to vanishing of the Bott--Chern group alone,
which gives the stated list.  In degree zero there is no nonzero
exact form, and in degree six integration identifies Bott--Chern and
de Rham cohomology.  This distinguishes the two meanings of a
``$\partial\bar\partial$-lemma in a degree''; a vanishing Bott--Chern
group in degree one does not make $\Delta^1$ vanish.
\end{proof}

\begin{claim}{C3.5}{Proved}{Metric obstructions that do not require $h$}
The threefold $X$ admits no balanced, strongly Gauduchon, K\"ahler,
locally conformally K\"ahler, or Hermitian-symplectic metric.
\end{claim}
\begin{proof}
Let $\theta=f^*\omega_{FS}$.  It is a nonzero semipositive exact
real $(1,1)$ form.  If $\omega$ were balanced, the real four-form
$\omega^2$ would be closed and
$\int_X\theta\wedge\omega^2>0$, because the integrand is positive
on the regular locus.  Exactness of $\theta$ and Stokes' theorem
give zero, a contradiction.

More generally, if $\omega$ were strongly Gauduchon, write
$\partial\omega^2=\bar\partial u$ with $u\in A^{3,1}$.
The real four-form $\Omega=\omega^2-u-\bar u$ is closed.
Bidegrees give $\theta\wedge\Omega=\theta\wedge\omega^2$,
so the same argument applies.  K\"ahler metrics are balanced.
If $d\omega=\ell\wedge\omega$ with a closed real Lee form
$\ell$, then $b_1=0$ gives $\ell=dv$, and
$d(e^{-v}\omega)=0$, excluding locally conformally K\"ahler metrics.
Finally any symplectic two-form on $X$ is exact because $b_2=0$.
Its cube has nonzero integral for its symplectic orientation, whereas
an exact cube has zero integral.  Thus $X$ is not symplectic and in
particular admits no Hermitian-symplectic metric.
\end{proof}

\begin{claim}{C3.6}{Known}{Gauduchon metrics with specified Chern--Ricci form}
There exists a Gauduchon metric on $X$.  More precisely, for every
Gauduchon metric $\omega_0$ there is another Gauduchon metric
$\omega$ with $[\omega^2]_A=[\omega_0^2]_A$ and
\begin{equation}\label{eq:C-Ricci}
\operatorname{Ric}^{\mathrm{Ch}}(\omega)=\pi f^*\omega_{FS}.
\end{equation}
This Ricci form is semipositive and has rank one at regular points
of $f$.  No Chern--Ricci-flat Hermitian metric exists on $X$.
\end{claim}
\begin{proof}
Gauduchon's conformal theorem
\citep[Theorem 2, pp. 132--133]{Gauduchon1985} applies to the compact
connected complex manifold $X$.  In complex dimension three, it
produces a positive conformal multiple $\omega_0$ of any Hermitian
form with $\partial\bar\partial\omega_0^2=0$.
For clarity, its reduction is a scalar elliptic equation: if
$L\phi$ is a nonzero fixed normalization of
$\operatorname{tr}_\omega(i\partial\bar\partial\phi)$, then
$L$ has kernel the constants by the maximum principle and Fredholm
index zero.  The one-dimensional kernel of $L^*$ has a strictly
positive generator $u$.  Indeed a sign-changing generator would
admit a nonzero nonnegative function orthogonal to it, placing such
a function in $\im L$, contrary to the maximum principle; the strong
maximum principle then gives strict positivity.  The form
$u^{1/2}\omega$ is Gauduchon.

The canonical-bundle formula in Part A and $4S_2=f^*(p_2)$ give
$K_X^2=f^*\OO(-1)$.  Consequently
$c_{1,BC}(T_X)=\tfrac12[\theta]$ with
$\theta=f^*\omega_{FS}$.  We use the conventional unnormalized
Chern--Ricci form $-i\partial\bar\partial\log\det g$,
which represents $2\pi c_{1,BC}(T_X)$.  Thus the real closed form
$\pi\theta$ represents its class.  Applying
\citet[Theorem 1.3]{STW2017} with background $\omega_0$ and prescribed
form $\pi\theta$ proves \eqref{eq:C-Ricci} and preservation of the
Aeppli class.  Finally C3.1 shows $[\theta]_{BC}\ne0$.  A metric
with zero Chern--Ricci form would make $c_{1,BC}(T_X)=0$, which is
impossible.
\end{proof}

\begin{claim}{C3.7}{Proved}{The fibre obstruction to SKT and astheno-K\"ahler metrics}
The threefold $X$ admits no pluriclosed (SKT) metric and no
astheno-K\"ahler metric.  The geometric input is the value $h=1$
and the global Dolbeault class $q$ in A3.2, whose restriction to a
regular fibre is $\eta=u\wedge w+6\gamma\wedge\delta$.
More generally, it is enough that the real image of
$H_A^{1,1}(X;\R)\to H_A^{1,1}(F;\R)$ be contained in
$\R\eta$ for one regular fibre $F$.
\end{claim}
\begin{proof}
We first compute the sign that matters.  Put
$y=\operatorname{Im}\tau>0$, $m=\operatorname{Im}\mu$, and
$\Delta=\operatorname{Im}\beta-6m^2/y<0$.
For the marked period rows
$\sigma_1=6\mu\gamma+\tau u+w$ and
$\sigma_2=\beta\gamma+\mu u+\delta$ we have
$\eta=u\wedge\sigma_1+6\gamma\wedge\sigma_2$.
This proves that the real class $\eta$ has type $(1,1)$.
The determinant of the four real period columns, with rows ordered
as $\operatorname{Re}\zeta_1,\operatorname{Im}\zeta_1,
\operatorname{Re}\zeta_2,\operatorname{Im}\zeta_2$, is
$y\Delta<0$.  Since $\eta^2=12\gamma\wedge u\wedge w\wedge\delta$,
\begin{equation}\label{eq:C-eta-negative}
\int_{F,\mathrm{or}_{\C}}\eta^2=-12.
\end{equation}
Equivalently the invariant representative of $\eta$ has Hermitian
coefficient matrix
\[
\frac1{y\Delta}
\begin{pmatrix}\operatorname{Im}\beta&-6m\\-6m&6y\end{pmatrix},
\qquad \det=\frac6{y\Delta}<0,
\]
so it has one positive and one negative eigenvalue.  The exact
symbolic determinant and an exact rational example are checked in
\texttt{code/c\_sign\_and\_fibre.py} by both eigenvalues and
diagonal elimination.

The explicit invariant fibre representative of $\gamma^{0,1}$ is
\[
\Gamma=\frac{i}{2\Delta}
\left(d\bar\zeta_2-\frac myd\bar\zeta_1\right).
\]
Writing $Q_F$ for the invariant $(1,1)$ representative of $\eta$,
direct exterior multiplication gives
\[
Q_F\wedge\Gamma=
\frac1{4y\Delta}
d\zeta_1\wedge d\bar\zeta_1\wedge d\bar\zeta_2,
\quad
\int_FQ_F\wedge\Gamma\wedge d\zeta_2=-1.
\]
This also proves the nonvanishing used in C1.2.

When $h=1$, C2.1 identifies $H_A^{1,1}(X)$ with the sum of the
one horizontal source and its vertical conjugate.  A representative
$Q$ of $q$ is nonzero in the horizontal source, since its fibre
restriction is nonzero whereas $d_1g$ restricts to zero.  Thus
$[Q]_A,[\bar Q]_A$ form a basis.  Their restrictions to $F$ both
equal the real class $\eta$.  The real image of the restriction map
is therefore $\R\eta$.

If $\omega$ were a pluriclosed positive real $(1,1)$ form, its
restriction to $F$ would be pluriclosed and positive.  Average this
restriction under translations of the compact complex torus $F$.
The average $\omega_{\mathrm{av}}$ is invariant, closed, and
strictly positive.  Translation preserves the Aeppli class: on a
torus the $\partial\bar\partial$-lemma identifies Aeppli
cohomology with the invariant $(p,q)$ forms, so the average is the
unique invariant representative of that class.  This identification
can also be checked Fourier mode by Fourier mode: each nonzero mode
is a Koszul square for two independent type components of its real
nonzero covector and contributes zero Aeppli cohomology.  Thus
$[\omega_{\mathrm{av}}]=r\eta$ for some real $r$.
But positivity gives $\int_F\omega_{\mathrm{av}}^2>0$, whereas
\eqref{eq:C-eta-negative} gives $-12r^2\le0$, a contradiction.
In dimension three the astheno-K\"ahler equation
$\partial\bar\partial\omega^{n-2}=0$ is exactly the SKT equation.
\end{proof}

\begin{claim}{C3.8}{Open}{A literal global coefficient atlas for the generators}
The invariant computations and all the metric types named above are
settled for $X$.  The real restriction map
\[
r_b:H_A^{1,1}(X;\R)\longrightarrow H_A^{1,1}(F_b;\R)
\]
has image $\R\eta$ by A3.2 and C3.7.  Its source has real dimension
two and its target has real dimension four.  This concerns
cohomology classes and does not assert literal equality of arbitrary
smooth representatives with an invariant fibre form.

The formula \eqref{eq:C-G-Green} supplies a global $G$ after a metric
choice.  For the standard primitive marking, $G_\gamma=G_0/12$
represents $g_\gamma$ by C3.2.  There is a precise operator
construction of $Q$ as well.  Fix a Hermitian metric on $X$ and a
regular fibre $F$, identify $\eta$ with its real invariant form on
$F$, and choose an orthonormal basis $e_1,e_2$ of the Dolbeault-harmonic
$(1,1)$ forms for $\Delta_{\bar\partial}$.  Define
\[
\ell(e)=-\frac1{12}\int_F i_F^*e\wedge\eta,\qquad
Q=\frac{\overline{\ell(e_1)}e_1+
            \overline{\ell(e_2)}e_2}
          {|\ell(e_1)|^2+|\ell(e_2)|^2}.
\]
The denominator is positive by A3.2 and
$\int_F\eta^2=-12$.  The expression is independent of the
orthonormal basis, is $\bar\partial$-closed, and has fibre
Dolbeault class $\eta$.  Indeed the image of fibre restriction on
$H^{1,1}_{\bar\partial}(X)$ is $\C\eta$, with kernel
$\C d_1g$; the displayed functional measures exactly the coefficient
of $\eta$, and the formula makes that coefficient one.

These two formulas fix global smooth representatives by explicit
elliptic operators and a finite normalization.  We have not evaluated
their coefficient functions on a common atlas through the three
special fibres.  The remaining task is therefore precise: compute
the harmonic projector and the two displayed restriction integrals
for one fully specified global Hermitian metric, or construct and
evaluate a Cech representative and partition-of-unity correction
with the same restrictions.  No numerical computation of those
infinite-dimensional elliptic inverses was run.  This is the
unfulfilled literal-coefficient requirement, not an undecided Hodge
number, product, or existence of the Leray class.
\end{claim}
